%% ED26-P02-syllogistic-0001. Exact source: arXiv:0808.0521v1, Sections 2--3.
\chapter{Categorical Syllogistic Inference}
\label{sec:p02-syllogistic-inference}

Categorical inference relates classes: inclusion, exclusion and the existence
of a shared member. The two calculi below give complete direct methods for
specified fragments of this reasoning. Their universe is nonempty, but a
named class may be empty. Consequently, a universal premise alone does not
usually supply an existential conclusion.

The presentation reconstructs Pratt-Hartmann and Moss's 2008 preprint,
including its two rule systems and the completeness argument of Section~3
\cite[pp.~3--18]{PrattHartmannMoss2008SyllogisticV1}.
These are modern formal systems, not Aristotle's. Boger's reconstruction of
Aristotle's underlying logic (Chapter~\ref{sec:p13-src063-0103}) writes the
predicate first, provides no term complement, makes every affirmation about a
nonexistent subject false and reads predication intensionally; Woods and
Irvine's reconstruction of Aristotle's early logic
\cite[pp.~54--56, 64--65]{WoodsIrvine2004EarlyLogic} restricts syllogisms by
premise minimality and non-circularity and excludes inconsistent premises.
Chapter~\ref{sec:p02-w02-p02-001-historical-contrast} states these
differences with exact sources and proves what each changes.

\section{Atoms, literals and the two fragments}

\begin{hbdefinition}{def:p02-syllogistic-language}{The languages \(\mathcal S\) and \(\mathcal S^\dagger\)}
Let \(P\) be a countably infinite set of unary predicate symbols, called
\emph{atoms}. Letters \(p,q,o\) range over atoms. A \emph{literal} is either
\(p\) or \(\bar p\), where \(p\in P\); letters \(l,m,n\) range over literals.
The bar on literals is an involution: \(\overline{\bar p}=p\).
A positive literal is an atom, and a negative literal is the complement of
an atom.

The formulas of \(\mathcal S\) are exactly
\[
 \exists(p,l),\quad \exists(l,p),\quad
 \forall(p,l),\quad \forall(l,\bar p).
\]
The formulas of \(\mathcal S^\dagger\) are exactly
\[
 \exists(l,m),\qquad \forall(l,m).
\]
Here \(\forall\) and \(\exists\) form complete categorical sentences from
two literals. They do not bind a written object-language variable. Arbitrary
conjunctions, disjunctions, Boolean combinations of terms, nested quantifiers
and binary relation symbols are not formulas of either fragment.
\end{hbdefinition}

Thus \(\forall(p,q)\) says that every \(p\) is a \(q\),
\(\forall(p,\bar q)\) that no \(p\) is a \(q\),
\(\exists(p,q)\) that some \(p\) is a \(q\), and
\(\exists(p,\bar q)\) that some \(p\) is not a \(q\).
The larger fragment additionally permits the genuinely new forms
\(\forall(\bar p,q)\) and \(\exists(\bar p,\bar q)\).
The smaller fragment already contains \(\forall(\bar p,\bar q)\) and
\(\exists(\bar p,q)\); the identifications below explain their role.
\index{syllogistic logic}\index{categorical inference}
\index{noun negation}

\begin{hbdefinition}{def:p02-syllogistic-semantics}{Set semantics and consequence}
A structure \(\mathfrak A\) consists of a nonempty universe \(A\) and a subset
\(p^{\mathfrak A}\subseteq A\) for each atom \(p\). Put
\((\bar p)^{\mathfrak A}=A\setminus p^{\mathfrak A}\). Define
\[
 \mathfrak A\models\forall(l,m)
 \quad\Longleftrightarrow\quad l^{\mathfrak A}\subseteq m^{\mathfrak A},
 \qquad
 \mathfrak A\models\exists(l,m)
 \quad\Longleftrightarrow\quad
 l^{\mathfrak A}\cap m^{\mathfrak A}\ne\varnothing.
\]
A set \(\Theta\) is satisfiable if one structure satisfies every member.
Write \(\Theta\models\theta\) when every structure satisfying \(\Theta\)
satisfies \(\theta\). No predicate is required to have a member.
\end{hbdefinition}

\begin{hbdefinition}{def:p02-syllogistic-identifications}{Formula complements and identifications}
Define
\[
 \overline{\forall(l,m)}=\exists(l,\bar m),\qquad
 \overline{\exists(l,m)}=\forall(l,\bar m).
\]
Each fragment is closed under this operation. In derivations identify
\(\exists(l,m)\) with \(\exists(m,l)\), and
\(\forall(l,m)\) with \(\forall(\bar m,\bar l)\).
Only these formula identifications are implicit. An \emph{absurdity} is a
formula \(\exists(l,\bar l)\) belonging to the selected fragment.
\end{hbdefinition}

\begin{hblemma}{thm:p02-syllogistic-semantic-identifications}{proved=annex:p2.syllogistic.identifications}{Meaning of the bar and formula identifications}
For every structure and formula \(\varphi\), exactly one of
\(\varphi,\bar\varphi\) is true. Both identifications preserve truth and
fragment membership; formula complementation respects them. Every absurdity
is false. For literals \(l,m\), let \(L(x),M(x)\) be their first-order
translations, using \(p(x)\) for \(p\) and \(\neg p(x)\) for \(\bar p\).
The sentences \(\forall(l,m)\) and \(\exists(l,m)\) have the same meanings
as, respectively,
\(\forall x(L(x)\to M(x))\) and
\(\exists x(L(x)\land M(x))\).
\end{hblemma}

\section{Direct rule systems}

\begin{hbdefinition}{def:p02-syllogistic-derivation}{Direct syllogistic derivation}
A substitution sends each atom to an atom and commutes with bars and the two
sentence constructors. It does not substitute arbitrary formulas or arbitrary
literals for atoms. Literal metavariables in the sheets abbreviate the
finitely many positive/negative rule patterns.

A direct derivation from \(\Theta\) is a finite tree. A leaf is a premise in
\(\Theta\) or an instance of a rule with no antecedents. Every other node is
the conclusion of an instance of a displayed rule whose antecedents occur
at its immediate predecessors, using the formula identifications already
declared. Write \(\Theta\vdash_{\mathsf S}\theta\) or
\(\Theta\vdash_{\mathsf S^\dagger}\theta\) for the corresponding relation.
The rules never discharge premises.
\end{hbdefinition}

\showcalculus{p02-syllogistic-s}
\showcalculus{p02-syllogistic-sdagger}

Rule A expresses empty-subject inclusion, T supplies reflexive inclusion,
and I retains a witness while forgetting one property. Rule X uses two
contradictory premises to infer a conclusion; it is not reductio ad absurdum.
Reductio would discharge a temporary assumption of \(\bar\theta\) after
deriving an absurdity. Neither direct calculus needs that additional step.
The extra rule N uses nonemptiness of the universe: if every non-\(l\)
is an \(l\), there are no non-\(l\)'s, so the whole universe consists of
\(l\)'s and supplies a witness.

\begin{hblemma}{thm:p02-syllogistic-rules}{proved=annex:p2.syllogistic.rules}{Sound rules and restriction to the smaller fragment}
Every rule in the two sheets is sound. A rule instance of
\(\mathsf S^\dagger\) whose antecedents and conclusion all belong to
\(\mathcal S\) is, under the stated formula identifications, an instance
of a rule of \(\mathsf S\); every \(\mathsf S\) rule instance is an
\(\mathsf S^\dagger\) rule instance under those identifications.
Rule D1 is redundant in \(\mathsf S\), though retained in its source sheet.
\end{hblemma}

\section{Separating universal information from witnesses}

The completeness construction exploits the unary language. A literal
specifies a class, a universal sentence constrains allowed combinations of
properties, and an existential sentence requests a witness to a combination.
Witnesses for different requests can be placed in disjoint components.

\begin{hbdefinition}{def:p02-syllogistic-disjoint-union}{Disjoint unions and restrictions}
For structures with pairwise disjoint universes \(A_i\), their disjoint union
has universe \(\bigcup_i A_i\) and atom interpretation
\(p^{\mathfrak A}=\bigcup_i p^{\mathfrak A_i}\). The index set must be
nonempty. A restriction to a nonempty subset \(B\subseteq A\) interprets
each atom by \(p^{\mathfrak A}\cap B\).
\end{hbdefinition}

\begin{hblemma}{thm:p02-syllogistic-separation}{proved=annex:p2.syllogistic.separation}{Universal constraints and existential premises}
Let \(\Phi\) contain only universal sentences and \(\Psi\) only existential
sentences, and suppose \(\Phi\cup\Psi\models\theta\).
\begin{enumerate}
\item If \(\Phi\cup\Psi\) is satisfiable and \(\theta\) is universal,
then \(\Phi\models\theta\).
\item If \(\Psi\ne\varnothing\) and \(\theta\) is existential, there is
\(\psi\in\Psi\) such that \(\Phi\cup\{\psi\}\models\theta\).
\item If \(\Psi=\varnothing\) and \(\theta=\exists(l,m)\), then
\(\Phi\models\forall(\bar l,l)\) and
\(\Phi\models\forall(\bar m,m)\).
\end{enumerate}
\end{hblemma}

\section{Literal closure and complete types}

\begin{hbdefinition}{def:p02-syllogistic-literal-closure}{Closure under universal derivations}
Write \(\Phi\vdash_{\mathrm{BTA}}\varphi\) for a derivation using only
the B, T and A rules of \(\mathsf S^\dagger\). For a set \(V\) of literals put
\[
 V^\Phi=\{m:\text{there is }l\in V
       \text{ with }\Phi\vdash_{\mathrm{BTA}}\forall(l,m)\}.
\]
A literal set is \emph{consistent} if it contains no pair \(l,\bar l\),
\emph{complete} if for every literal \(l\) it contains \(l\) or \(\bar l\),
and \(\Phi\)-\emph{closed} if \(V^\Phi\subseteq V\).
These literal-set properties differ from completeness of a proof system.
\end{hbdefinition}

\begin{hblemma}{thm:p02-syllogistic-closure}{proved=annex:p2.syllogistic.closure}{Closure and a finite conflict witness}
For all \(\Phi,V\), one has \(V\subseteq V^\Phi\), and \(V^\Phi\) is
\(\Phi\)-closed. A union of \(\Phi\)-closed literal sets is
\(\Phi\)-closed. If \(V^\Phi\) is inconsistent, some \(l,l'\in V\) satisfy
\(\Phi\vdash_{\mathrm{BTA}}\forall(l,\bar l')\).
\end{hblemma}

\begin{hblemma}{thm:p02-syllogistic-extension}{proved=annex:p2.syllogistic.extension}{Extending a consistent closed literal set}
Every nonempty, consistent, \(\Phi\)-closed literal set has a consistent,
complete, \(\Phi\)-closed extension. Such a complete extension defines a
one-element structure satisfying \(\Phi\), with its unique element satisfying
exactly the literals in the extension.
\end{hblemma}

The nonempty starting-set condition matters. A universally inconsistent
theory can have the empty set as a consistent closed literal set; the lemma
must not then assert the existence of a model.

\section{Completeness and a finite decision method}

\begin{hbtheorem}{thm:p02-syllogistic-sdagger-completeness}{proved=annex:p2.syllogistic.sdagger}{Soundness and completeness of \(\mathsf S^\dagger\)}
For every set \(\Theta\subseteq\mathcal S^\dagger\) and formula
\(\theta\in\mathcal S^\dagger\),
\[
 \Theta\models\theta
 \quad\Longleftrightarrow\quad
 \Theta\vdash_{\mathsf S^\dagger}\theta.
\]
The result includes inconsistent premise sets and uses no reductio rule.
\end{hbtheorem}

\begin{hbtheorem}{thm:p02-syllogistic-s-completeness}{proved=annex:p2.syllogistic.s}{Soundness, completeness and conservativity for \(\mathsf S\)}
For \(\Theta\subseteq\mathcal S\) and \(\theta\in\mathcal S\),
\[
 \Theta\models\theta
 \quad\Longleftrightarrow\quad
 \Theta\vdash_{\mathsf S}\theta
 \quad\Longleftrightarrow\quad
 \Theta\vdash_{\mathsf S^\dagger}\theta.
\]
Adding noun negation in the larger language creates new sentences, but no new
consequence between premises and conclusions of the smaller fragment.
\end{hbtheorem}

The proofs follow the source's one-element countermodel construction and its
restriction argument, including every satisfiable and inconsistent case
\cite[Section~3, Lemmas~3.1--3.5 and Theorems~3.6--3.7]
{PrattHartmannMoss2008SyllogisticV1}. For arbitrary premise sets, the use of
first-order compactness is explicit; it reduces the task to a finite subset
without replacing the direct syllogistic completeness argument.

\begin{hbmethod}{def:p02-syllogistic-saturation}{Finite saturation with proof records}
For a finite input \(\Theta,\theta\), let \(Q\) be its occurring atoms.
List all sentences of the chosen fragment using \(Q\), with one representative
for each declared identification. Begin with the premises and all T axioms.
Repeatedly add the conclusion of every rule instance over \(Q\) whose
antecedents have already been added. Store the rule and predecessor sentences
for each new conclusion. Stop when a full pass adds nothing.
Return entailment exactly when the closure contains \(\theta\); the stored
predecessors give a finite direct proof. An omitted conclusion means
non-entailment, and hence the existence of a countermodel, by completeness.
\end{hbmethod}

\begin{hbtheorem}{thm:p02-syllogistic-saturation}{proved=annex:p2.syllogistic.saturation}{Termination and correctness of saturation}
The method decides consequence in either fragment for finite inputs in
polynomial time. There are at most \(8|Q|^2\) sentence representatives before
identifications, and no fresh atom is needed in a derivation of an input
conclusion. This bound does not assert that the displayed implementation has
the best possible space or time complexity.
\end{hbtheorem}

\section{Worked inferences and existential import}

\begin{longtable}{@{}p{0.29\linewidth}p{0.25\linewidth}p{0.35\linewidth}@{}}
\caption{Categorical inference and its existence requirements}
\label{tab:p02-syllogistic-lookup}\\
\hline
Premises & Proposed conclusion & Condition or countermodel \\
\hline
\endfirsthead
\caption[]{Categorical inference and its existence requirements (continued)}\\
\hline
Premises & Proposed conclusion & Condition or countermodel \\
\hline
\endhead
\hline\multicolumn{3}{r}{\emph{continued on next page}}\\
\endfoot
\endlastfoot
\(\forall(p,q)\), \(\forall(q,o)\) & \(\forall(p,o)\) & Valid by B, even when \(p\) is empty. \\
\(\forall(p,q)\) & \(\exists(p,q)\) & Requires a witness such as \(\exists(p,p)\); an empty \(p\) refutes unrestricted existential import. \\
\(\forall(p,\bar q)\) & \(\forall(q,\bar p)\) & Valid by universal contraposition; no existence premise is needed. \\
\(\exists(p,q)\) & \(\exists(q,p)\) & Valid by existential symmetry; the same element is the witness. \\
\(\forall(p,q)\) & \(\forall(q,p)\) & Invalid: in a singleton universe, take \(p^{\mathfrak A}=\varnothing\) and \(q^{\mathfrak A}=A\). \\
\(\forall(\bar p,p)\) & \(\exists(p,p)\) & Valid in \(\mathsf S^\dagger\) by N; the antecedent forces \(p^{\mathfrak A}=A\ne\varnothing\). \\
\hline
\end{longtable}

\begin{hbexample}{def:p02-syllogistic-oaks}{Inclusion and an explicit witness}
Let \(o,t,p\) mean oak, tree and plant. From
\(\forall(o,t)\) and \(\forall(t,p)\), rule B yields \(\forall(o,p)\).
No existence premise is needed for this universal conclusion.
If the additional premise \(\exists(o,o)\) is supplied, D1 with
\(\forall(o,p)\) yields \(\exists(o,p)\). The latter conclusion is warranted
because the premise supplied a witness; the two inclusions did not create it.
\end{hbexample}

\begin{hbexample}{def:p02-syllogistic-negative-witness}{Passing a witness through a positive and a negative inclusion}
Let \(a,b,c,d\) mean artist, beekeeper, carpenter and dentist. Use premises
\(\exists(a,b)\), \(\forall(a,c)\) and \(\forall(b,\bar d)\).
D2 first gives \(\exists(c,b)\). D1 then gives
\(\exists(c,\bar d)\): some carpenter is not a dentist. This is the source's
unary worked pattern, with every applied rule remaining in \(\mathcal S\).
\end{hbexample}

\begin{hbexample}{def:p02-syllogistic-empty-term}{An empty subject blocks existential import}
Let \(A=\{s\}\), with \(s\) a stone, and interpret oak, tree and plant by
the empty set. Both oak inclusions in
Example~\ref{def:p02-syllogistic-oaks} are true, but
\(\exists(o,p)\) is false. Thus the universe's nonemptiness does not make
any particular predicate nonempty. Likewise, \(\forall(o,o)\) is valid while
\(\exists(o,o)\) is not.
\end{hbexample}

\begin{hbexample}{def:p02-syllogistic-nonempty-universe}{The exact role of rule N}
In \(\mathcal S^\dagger\), \(\forall(\bar p,p)\) forces \(p\) to contain
the whole universe; N therefore yields \(\exists(p,p)\). If an empty universe
were permitted, the antecedent would be true there and the conclusion false.
Conversely, \(\forall(p,\bar p)\) forces \(p\) empty; N applied with
\(l=\bar p\) gives \(\exists(\bar p,\bar p)\). This new negative existential
is not an \(\mathcal S\)-formula.
\end{hbexample}

\begin{hbexample}{def:p02-syllogistic-inconsistent-premises}{Inconsistent categorical premises}
From \(\exists(p,p)\) and \(\forall(p,\bar p)\), the two premises are
formula complements. X yields any sentence of the selected fragment. This
is the same vacuity expressed by semantic consequence from an unsatisfiable
premise set; it supplies no evidence that the premises describe a possible
population.
\end{hbexample}

\begin{hbconvention}{def:p02-syllogistic-lookup}{Task and scope lookup}
Use \(\mathsf S\) for the four categorical forms with positive subjects after
the declared identifications. Use \(\mathsf S^\dagger\) when a genuinely
negative subject or a witness satisfying two negative literals is required.
For a requested existential conclusion, first locate the witness premise,
or a universal constraint forcing the relevant classes to fill a nonempty
universe. For a failed inference, retain empty extensions when searching for
a countermodel. Relational predicates, arbitrary Boolean terms and the
source's relational fragments require their own calculi; the present theorems
do not cover them.
\end{hbconvention}

Aristotle's fourteen two-premise syllogism rules, in Boger's reconstruction
(Chapter~\ref{sec:p13-src063-0103}), assume existent subjects.
Chapter~\ref{sec:p02-w02-p02-002-fourteen} reads them in \(\mathcal S\):
twelve need no existence premise, while Darapti and Felapton need the
premise that their middle term is nonempty and fail in a structure where it
is empty.

\index{existential import}\index{syllogistic logic!completeness}
\index{syllogistic logic!empty terms}\index{syllogistic logic!direct derivation}
